% Manuscript fragment. Requires amsmath, amssymb, amsthm.
\subsection{Explicit guarded normalization of record readability}
\label{sec:guarded-provenance}

One session version can contain several reserved values, and several
records can contain the same value. Suppose we are given the objects
and keys acquired up to time $t$. Which original draws can we now read?
A source $o\in\Omega_{\rm src}$ names one such draw, and
$\mathcal E^{\rm src}_t$ records the identities of the sources whose
protected values we can derive. We continue to use $\mathcal C_t$ and
$\mathcal L_t$ for versions and $\mathcal O[v]$ for the payload of $v$.

Here we fix the graph of admitted records and the public key-derivation
edges. Key derivation in this fragment is independent of recovered
plaintext. The preceding interface specifies how to add an explicit
plaintext-dependent key-extraction rule. The provenance rules supply
the identities in the graph, and the renderer supplies the corresponding
bytes. A second copy carries the same source identity and can give the
adversary another route to its value.

\paragraph{Immutable record graph.}
Let $\mathcal R$ be a finite set of immutable encrypted record objects.
Each $r\in\mathcal R$ has a matching decryption-key identity $\kappa(r)$,
a set of embedded child records $\operatorname{child}(r)$, and a set
$\operatorname{orig}(r)$ of protected source values in its plaintext.
Containment is acyclic.  Record identities distinguish ciphertext objects;
source identities distinguish original sampling operations.  Several
records can contain the same source.  A rewrapped record has a new record
identity and may have a new key, while retaining its source identity.
Key identities are equivalence classes of equal actual keys. Equal key bytes
therefore have one identity across separate API objects.

Let $\mathcal D\subseteq\mathcal K\times\mathcal K$ be fixed public key
derivation edges.  An edge $(k,k')$ means that possession of $k$ permits
deterministic computation of $k'$.  This interface represents, for example,
derivation of an operation-wrapping key from the journal key.  In the
theorem below, these edges do not depend on plaintext recovered from records.
Write $D^*$ for their reflexive transitive closure.

Let $\Omega_{\rm src}=\bigcup_{r\in\mathcal R}\operatorname{orig}(r)$ be the
declared source domain.  At a trace prefix $t$, let $A_t\subseteq\mathcal R$ be all record objects
ever acquired by that prefix, $K_t\subseteq\mathcal K$ all directly
disclosed keys, and $\mathcal S^{\rm dir}_t\subseteq\Omega_{\rm src}$ all directly disclosed source values, represented
by their source identities.  These are cumulative sets.  The decryption-key
closure is
\[
  K_t^*=\{k':\exists k\in K_t,\ (k,k')\in D^*\}.
\]
The concrete derivation rules are
\[
 \frac{k\quad(k,k')\in\mathcal D}{k'},\qquad
 \frac{r\quad\kappa(r)}{c}\ (c\in\operatorname{child}(r)),\qquad
 \frac{r\quad\kappa(r)}{o}\ (o\in\operatorname{orig}(r)).
\]
Their least closure starts from $A_t,K_t,\mathcal S^{\rm dir}_t$.  A fact $o$ in the source
sort means that the protected value of source $o$ is derivable; it does
not mean that a public identifier has become known.  The source projection
of this closure is denoted $\mathcal E^{\rm src}_t$.

\paragraph{Guarded routes.}
For every path $p=(r_0,\ldots,r_\ell)$ in the containment graph ending
at a record with $o\in\operatorname{orig}(r_\ell)$, define
\[
 g(p)=\{\kappa(r_0),\ldots,\kappa(r_\ell)\}.
\]
The normalized table contains a route $(r_0,g(p),o,p)$ for every such
path.  Its exposure expression is
\[
 \widehat{\mathcal E}^{\rm src}_t=\mathcal S^{\rm dir}_t\ \cup\
 \{o:\exists(r,g,o,p),\ r\in A_t\ \land\ g\subseteq K_t^*\}.
 \tag{GP}
\]
One available path to $o$ suffices. Along that path every required
key must be available, and a key used twice need only be known once.

\begin{theorem}[Exact guarded normalization]
\label{thm:guarded-normalization}
For every finite record graph and derivation interface above, and every
acquisition prefix $(A_t,K_t,\mathcal S^{\rm dir}_t)$, the concrete and normalized exposure
sets agree: $\mathcal E^{\rm src}_t=\widehat{\mathcal E}^{\rm src}_t$.  The equality applies to all prefixes of
an adaptively chosen history.
\end{theorem}
\begin{proof}
Key derivations depend only on disclosed keys and the fixed relation
$\mathcal D$, so their least closure is exactly $K_t^*$.
For the forward inclusion, take a source $o\in \mathcal E^{\rm src}_t$.
If it was directly disclosed then $o\in \mathcal S^{\rm dir}_t$.
Otherwise its first derivation decrypts a record $r_\ell$ with the
matching key.  If that record was not acquired directly, its first
availability follows by decrypting a parent.  Follow these first
derivations backwards.  The finite acyclic record graph makes this
process terminate at an acquired record $r_0$.  Every intervening key
lies in $K_t^*$, yielding a route satisfying (GP).
For the reverse inclusion, take a satisfied route.  Starting with its
acquired first record, successively decrypt along the path.  Every
required key is in $K_t^*$, so each child and finally $o$ is in the
concrete least closure.  Direct disclosures are in that closure by
definition.  Both inclusions hold for each fixed prefix; no independence
or advance choice of the acquisition schedule was used.
\end{proof}

\paragraph{Incremental evaluation.}
Equation (GP) answers the readability question by listing the paths to
each source. A small record DAG can have exponentially many such paths.
As acquisitions and disclosures arrive, we can instead track the facts
that each event makes derivable. The evaluator uses indexed forward
chaining for Horn rules~\cite[Section~4]{dowlingGallier1984}.
Write
\[
\begin{aligned}
 N&=|\mathcal R|+|\mathcal K|+|\Omega_{\rm src}|,\\
 M&=|\mathcal D|+\sum_{r\in\mathcal R}
       \bigl(1+|\operatorname{child}(r)|+|\operatorname{orig}(r)|\bigr).
\end{aligned}
\]
The unit in each summand counts the record's matching-key incidence.
In an adjacency-list representation, repeated incidences are counted with
multiplicity in $M$.

\begin{proposition}[Incremental exposure evaluation]
\label{prop:guarded-worklist}
For the fixed record and key-derivation graphs above, the exact exposure
set at every prefix of a cumulative acquisition and disclosure history
can be maintained with $O(N+M)$ preprocessing, $O(N+M+U)$ total time and
$O(N+M)$ working space, where $U$ counts input acquisition, key-disclosure
and direct-source-disclosure events, including repetitions. These bounds
use indexed atom identities with constant-time membership and insertion;
they include reporting each newly exposed source once. Exporting complete
sets at selected prefixes adds their total output size.
\end{proposition}
\begin{proof}
Index each key's outgoing derivation edges and the records having that
matching key. Maintain known keys, available records, opened records and
exposed sources. Enqueue a key or record only when it first becomes known
or available. A key event propagates its outgoing derivations and checks
its indexed records. A record event checks its matching key. When both
the record and its key are present, open the record once, making its
children available and its sources exposed. Direct disclosures insert
their sources immediately. Drain the queue after each input event.

Every insertion is an instance of a stated derivation rule, so all
maintained facts belong to the least closure. At an empty queue, every
derivation edge with a known parent has been processed. For a record
whose two premises hold, the later arrival of those premises triggers
its opening unless it was already opened. Thus every enabled rule has
fired, and the maintained facts are exactly the least closure. Applying
this argument after each input event and using
Theorem~\ref{thm:guarded-normalization} gives the exposure set (GP) at
every prefix, including adaptively chosen prefixes.

Each key and record is queued at most once. Each key-derivation and
matching-key incidence is scanned at most once from its key, each
record's key is tested once from its record, and each opened record's
payload incidences are processed once. Repeated input events take
constant time. The total work is therefore $O(N+M+U)$, and the indexes,
sets and queue occupy $O(N+M)$ space. Each source enters the incremental
output at most once; complete-set exports are charged separately.
\end{proof}

The graphs are fixed throughout this bound. Local deletion changes future
acquisition availability and leaves the cumulative closure intact. A new
record or derivation edge changes the indexed input and requires its own
update accounting. An implementation using dictionaries has the
corresponding expected bounds under constant expected lookup and insertion
costs; normalization of external identities is charged separately.
The finite checker retains explicit route enumeration
for its small fixtures, where it compares the guarded normal form with
forward decryption. The indexed evaluator avoids that enumeration;
materializing all route witnesses can itself require exponential output.

\begin{corollary}[Copies, deletion and late disclosure]
\label{cor:guarded-history}
Exposure is monotone under additional acquisitions, direct disclosures
and key disclosures. Copying or rewrapping a source preserves existing
routes and may add routes to the same source identity.
Local deletion of a record previously acquired
does not remove its routes from adversarial knowledge.  In particular,
a later matching-key disclosure can expose a source in a ciphertext
acquired before that disclosure and deleted locally in the meantime.
\end{corollary}
\begin{proof}
The cumulative sets and their key closure only grow, and every predicate
in (GP) is positive. An exact ciphertext copy reuses its record identity
and route. Source-preserving rewrapping may change record identities or
key guards while retaining the terminal source. Local deletion
restricts future acquisition from that local store; it removes no member
of $A_t$ and therefore cannot invalidate a satisfied route.
\end{proof}

\paragraph{What a readable source tells us.}
For a source $o$, a path from an acquired record and the keys that
open every wrapper on that path give a witness for its readability.
Direct disclosure enters $o$ in $\mathcal E^{\rm src}_t$ immediately.
Absence from the set has a precise meaning within this fragment. These
record decryption and fixed public key-derivation rules have no way
to derive the value from the history so far. Other ways of learning a
secret, including guessing, cryptanalysis, timing leakage or access to
a quantum decryption oracle, require their own model. For a claim
about the selected application message we simulate the native
ciphertexts and commitments and bound the adversary's advantage in
the protocol experiment.

\paragraph{Native shape and finite checks.}
At the fixed research source, operation reservations are encrypted under
a wrapping key derived from the journal key, and tokens are retained
within encrypted durable images or pending targets.  The corresponding
route therefore includes the image key and the token key, with their
actual derivation edge.  A fresh aggregate image can preserve an older
token and its original source.  Public image-head digests are distinct
observations from the record objects needed by the decryption rules.
The executable supplement independently evaluates iterative decryption
and guarded-route normalization, exhausts all acquisition/key/direct-source
subsets for its declared finite graphs, and checks histories with late
key disclosure and local deletion.  Its counts and source hash are in
the accompanying machine-readable result; those checks validate the
implemented finite model and supplement the general proof above.
